% Options for packages loaded elsewhere
\PassOptionsToPackage{unicode}{hyperref}
\PassOptionsToPackage{hyphens}{url}
\documentclass[
  ignorenonframetext,
]{beamer}
\newif\ifbibliography
\usepackage{pgfpages}
\setbeamertemplate{caption}[numbered]
\setbeamertemplate{caption label separator}{: }
\setbeamercolor{caption name}{fg=normal text.fg}
\beamertemplatenavigationsymbolsempty
% remove section numbering
\setbeamertemplate{part page}{
  \centering
  \begin{beamercolorbox}[sep=16pt,center]{part title}
    \usebeamerfont{part title}\insertpart\par
  \end{beamercolorbox}
}
\setbeamertemplate{section page}{
  \centering
  \begin{beamercolorbox}[sep=12pt,center]{section title}
    \usebeamerfont{section title}\insertsection\par
  \end{beamercolorbox}
}
\setbeamertemplate{subsection page}{
  \centering
  \begin{beamercolorbox}[sep=8pt,center]{subsection title}
    \usebeamerfont{subsection title}\insertsubsection\par
  \end{beamercolorbox}
}
% Prevent slide breaks in the middle of a paragraph
\widowpenalties 1 10000
\raggedbottom
\AtBeginPart{
  \frame{\partpage}
}
\AtBeginSection{
  \ifbibliography
  \else
    \frame{\sectionpage}
  \fi
}
\AtBeginSubsection{
  \frame{\subsectionpage}
}
\usepackage{iftex}
\ifPDFTeX
  \usepackage[T1]{fontenc}
  \usepackage[utf8]{inputenc}
  \usepackage{textcomp} % provide euro and other symbols
\else % if luatex or xetex
  \usepackage{unicode-math} % this also loads fontspec
  \defaultfontfeatures{Scale=MatchLowercase}
  \defaultfontfeatures[\rmfamily]{Ligatures=TeX,Scale=1}
\fi
\usepackage{lmodern}
\ifPDFTeX\else
  % xetex/luatex font selection
\fi
% Use upquote if available, for straight quotes in verbatim environments
\IfFileExists{upquote.sty}{\usepackage{upquote}}{}
\IfFileExists{microtype.sty}{% use microtype if available
  \usepackage[]{microtype}
  \UseMicrotypeSet[protrusion]{basicmath} % disable protrusion for tt fonts
}{}
\makeatletter
\@ifundefined{KOMAClassName}{% if non-KOMA class
  \IfFileExists{parskip.sty}{%
    \usepackage{parskip}
  }{% else
    \setlength{\parindent}{0pt}
    \setlength{\parskip}{6pt plus 2pt minus 1pt}}
}{% if KOMA class
  \KOMAoptions{parskip=half}}
\makeatother
\usepackage{color}
\usepackage{fancyvrb}
\newcommand{\VerbBar}{|}
\newcommand{\VERB}{\Verb[commandchars=\\\{\}]}
\DefineVerbatimEnvironment{Highlighting}{Verbatim}{commandchars=\\\{\}}
% Add ',fontsize=\small' for more characters per line
\usepackage{framed}
\definecolor{shadecolor}{RGB}{248,248,248}
\newenvironment{Shaded}{\begin{snugshade}}{\end{snugshade}}
\newcommand{\AlertTok}[1]{\textcolor[rgb]{0.94,0.16,0.16}{#1}}
\newcommand{\AnnotationTok}[1]{\textcolor[rgb]{0.56,0.35,0.01}{\textbf{\textit{#1}}}}
\newcommand{\AttributeTok}[1]{\textcolor[rgb]{0.13,0.29,0.53}{#1}}
\newcommand{\BaseNTok}[1]{\textcolor[rgb]{0.00,0.00,0.81}{#1}}
\newcommand{\BuiltInTok}[1]{#1}
\newcommand{\CharTok}[1]{\textcolor[rgb]{0.31,0.60,0.02}{#1}}
\newcommand{\CommentTok}[1]{\textcolor[rgb]{0.56,0.35,0.01}{\textit{#1}}}
\newcommand{\CommentVarTok}[1]{\textcolor[rgb]{0.56,0.35,0.01}{\textbf{\textit{#1}}}}
\newcommand{\ConstantTok}[1]{\textcolor[rgb]{0.56,0.35,0.01}{#1}}
\newcommand{\ControlFlowTok}[1]{\textcolor[rgb]{0.13,0.29,0.53}{\textbf{#1}}}
\newcommand{\DataTypeTok}[1]{\textcolor[rgb]{0.13,0.29,0.53}{#1}}
\newcommand{\DecValTok}[1]{\textcolor[rgb]{0.00,0.00,0.81}{#1}}
\newcommand{\DocumentationTok}[1]{\textcolor[rgb]{0.56,0.35,0.01}{\textbf{\textit{#1}}}}
\newcommand{\ErrorTok}[1]{\textcolor[rgb]{0.64,0.00,0.00}{\textbf{#1}}}
\newcommand{\ExtensionTok}[1]{#1}
\newcommand{\FloatTok}[1]{\textcolor[rgb]{0.00,0.00,0.81}{#1}}
\newcommand{\FunctionTok}[1]{\textcolor[rgb]{0.13,0.29,0.53}{\textbf{#1}}}
\newcommand{\ImportTok}[1]{#1}
\newcommand{\InformationTok}[1]{\textcolor[rgb]{0.56,0.35,0.01}{\textbf{\textit{#1}}}}
\newcommand{\KeywordTok}[1]{\textcolor[rgb]{0.13,0.29,0.53}{\textbf{#1}}}
\newcommand{\NormalTok}[1]{#1}
\newcommand{\OperatorTok}[1]{\textcolor[rgb]{0.81,0.36,0.00}{\textbf{#1}}}
\newcommand{\OtherTok}[1]{\textcolor[rgb]{0.56,0.35,0.01}{#1}}
\newcommand{\PreprocessorTok}[1]{\textcolor[rgb]{0.56,0.35,0.01}{\textit{#1}}}
\newcommand{\RegionMarkerTok}[1]{#1}
\newcommand{\SpecialCharTok}[1]{\textcolor[rgb]{0.81,0.36,0.00}{\textbf{#1}}}
\newcommand{\SpecialStringTok}[1]{\textcolor[rgb]{0.31,0.60,0.02}{#1}}
\newcommand{\StringTok}[1]{\textcolor[rgb]{0.31,0.60,0.02}{#1}}
\newcommand{\VariableTok}[1]{\textcolor[rgb]{0.00,0.00,0.00}{#1}}
\newcommand{\VerbatimStringTok}[1]{\textcolor[rgb]{0.31,0.60,0.02}{#1}}
\newcommand{\WarningTok}[1]{\textcolor[rgb]{0.56,0.35,0.01}{\textbf{\textit{#1}}}}
\usepackage{graphicx}
\makeatletter
\newsavebox\pandoc@box
\newcommand*\pandocbounded[1]{% scales image to fit in text height/width
  \sbox\pandoc@box{#1}%
  \Gscale@div\@tempa{\textheight}{\dimexpr\ht\pandoc@box+\dp\pandoc@box\relax}%
  \Gscale@div\@tempb{\linewidth}{\wd\pandoc@box}%
  \ifdim\@tempb\p@<\@tempa\p@\let\@tempa\@tempb\fi% select the smaller of both
  \ifdim\@tempa\p@<\p@\scalebox{\@tempa}{\usebox\pandoc@box}%
  \else\usebox{\pandoc@box}%
  \fi%
}
% Set default figure placement to htbp
\def\fps@figure{htbp}
\makeatother
\setlength{\emergencystretch}{3em} % prevent overfull lines
\providecommand{\tightlist}{%
  \setlength{\itemsep}{0pt}\setlength{\parskip}{0pt}}
\usepackage{bookmark}
\IfFileExists{xurl.sty}{\usepackage{xurl}}{} % add URL line breaks if available
\urlstyle{same}
\hypersetup{
  pdftitle={TMA4315 Generalized linear models H2018},
  pdfauthor={Mette Langaas and Ingeborg Hem, Department of Mathematical Sciences, NTNU},
  hidelinks,
  pdfcreator={LaTeX via pandoc}}

\title{TMA4315 Generalized linear models H2018}
\subtitle{Module 8: Generalized linear mixed effects models (GLMM)}
\author{Mette Langaas and Ingeborg Hem, Department of Mathematical
Sciences, NTNU}
\date{15.11.2018 {[}PL{]}, 16.11.2018 {[}IL{]}}

\begin{document}
\frame{\titlepage}

\begin{frame}
(Latest changes: 16.11: IL added links, 13.11.2018 - first version).
\end{frame}

\begin{frame}{Overview}
\phantomsection\label{overview}
Aim: Present methods for analysing correlated responses in a generalized
linear models setting - LMM meets GLM to become \emph{GLMM}.

Also here (as in module LMM) we will only consider \emph{two-level
models} and in particular focus on \emph{random intercept models} for
binomial and Poisson responses. Emphasis will be on understanding.
\end{frame}

\begin{frame}
\begin{block}{Learning material}
\phantomsection\label{learning-material}
\begin{itemize}
\tightlist
\item
  Textbook: Fahrmeir et al (2013): Chapter 7.5: 389-394, 7.7.
\item
  \href{https://www.math.ntnu.no/emner/TMA4315/2018h/M8PL.pdf}{Classnotes
  (15.11.2018)}
\end{itemize}
\end{block}
\end{frame}

\begin{frame}
\begin{block}{Topics}
\phantomsection\label{topics}
\begin{itemize}
\tightlist
\item
  beaches example - revisited
\item
  notation
\item
  the generalized linear mixed effect model (three ingredients)
\item
  the GLMM with random intercept
\item
  the marginal model
\item
  parameter estimation and Laplace approximation
\item
  summing up: what do we need to know about the GLMM?
\item
  \hyperlink{additional}{additional info on different software (not on
  reading list)}
\end{itemize}
\end{block}
\end{frame}

\begin{frame}{Beaches example - revisited}
\phantomsection\label{beaches-example---revisited}
This example is taken from Zuur et al.~(2009, chapter 5, pages 101-142),
and data are referred to as RIKZ. Data were collected on nine different
beaches in the Netherlands with the aim to investigate if there is a
relationship between:

\begin{itemize}
\tightlist
\item
  richness of species (number of species observed) and
\item
  NAP: the height of the sampling station compared to mean tidal level.
\end{itemize}

Data: 45 observations, taken at 9 beaches:

\begin{itemize}
\tightlist
\item
  beach: the beach that the samples were taken, for each beach 5
  different samples were taken.
\end{itemize}
\end{frame}

\begin{frame}
\pandocbounded{\includegraphics[keepaspectratio]{Module08GLMMPresentation_files/figure-beamer/unnamed-chunk-1-1.pdf}}
\end{frame}

\begin{frame}
We want to use the richness of species as the response in a regression
model.

In Module 7 we assumed that the response could be viewed as a normal
variable, but now we assume that the response comes from a \emph{Poisson
distribution}.

As in the LMM module, we could (1) ignore the beach effect, (2) include
it as a fixed effect, or (3) use a random effect. Here we will use (3),
and extend the GLM framework to include random effects
\end{frame}

\begin{frame}
But - first: what do we remember about GLMs (and Poission regression
with log-link in particular)?
\end{frame}

\begin{frame}
\textbf{Assumptions for GLM model (Poisson): }

\begin{enumerate}
\item
  Random component from exponential family: We have for our beaches:
  \(Y_i \sim \text{Poisson}(\lambda_i)\), with
  \(\text{E}(Y_i)=\lambda_i\), and \(\text{Var}(Y_i)=\lambda_i\).
\item
  Linear predictor: \(\eta_i={\bf x}_i^T {\boldsymbol \beta}\).
\item
  Log link \[\eta_i=\ln(\lambda_i)=g(\lambda_i)\] and (inverse thereof)
  response function \[\lambda_i=\exp(\eta_i)\] Canonical link for
  Poisson is log.
\end{enumerate}
\end{frame}

\begin{frame}{Generalized linear mixed effects models}
\phantomsection\label{generalized-linear-mixed-effects-models}
\begin{block}{Notation}
\phantomsection\label{notation}
combining GLM (Module 5) and LMM (Module 7)

We have \(n_i\) repeated observations (\(j=1,\ldots, n_i\)) from each of
\(i=1,\ldots,m\) clusters (or individuals).

\begin{itemize}
\tightlist
\item
  Responses: \(Y_{i1},Y_{i2},\ldots, Y_{in_i}\) (e.g.~species richness
  at beach \(i\) sample \(j\))
\item
  Covariates: \({\bf x}_{i1},{\bf x}_{i1}, \ldots, {\bf x}_{in_i}\)
  (e.g.~NAP for beach \(i\) sample \(j\))
\end{itemize}

The covariates \({\bf x}_{ij}\) are \(p\times 1\) vectors (as before,
\(k\) covariates and one intercept so \(p=k+1\)).
\end{block}
\end{frame}

\begin{frame}
\begin{block}{Repetition GLM}
\phantomsection\label{repetition-glm}
\begin{enumerate}
\tightlist
\item
  \textbf{Random component} - exponential family (now \(y_{ij}\))
\end{enumerate}

\[ f(y_{ij}\mid \theta_{ij})=\exp \left( \frac{y_{ij} \theta_{ij}-b(\theta_{ij})}{\phi}\cdot w_{ij} + c(y_{ij}, \phi, w_{ij}) \right) \]

\begin{itemize}
\item
  \(\theta_{ij}\) is called the canonical parameter and is a parameter
  of interest
\item
  \(\phi\) is called a nuisance parameter (and is not of interest to
  us=therefore a nuisance (plage))
\item
  \(w_{ij}\) is a weight function, in most cases \(w_{ij}=1\)
\item
  \(b\) and \(c\) are known functions.
\end{itemize}

We looked in detail into the binomial, Poisson and gamma distribution
(in addition to the normal).
\end{block}
\end{frame}

\begin{frame}
\begin{enumerate}
\setcounter{enumi}{1}
\tightlist
\item
  \textbf{Linear predictor} (previously ``systematic component''):
\end{enumerate}

\[\eta_{ij}={\bf x}_{ij}^T {\boldsymbol \beta}\]

\begin{enumerate}
\setcounter{enumi}{2}
\tightlist
\item
  \textbf{Link function \(g\)- and response function \(h\)}
\end{enumerate}

\[\eta_{ij}=g(\mu_{ij})\]

\[\mu_{ij}=h(\eta_{ij})\]
\end{frame}

\begin{frame}
\begin{block}{New element}
\phantomsection\label{new-element}
we add a random component (effect) to the \emph{linear predictor}!

\begin{enumerate}
\setcounter{enumi}{1}
\tightlist
\item
  \[\eta_{ij}={\bf x}_{ij}^T {\boldsymbol \beta}+ {\bf u}^T_{ij} \gamma_i\]
  and we make the same assumptions as for the LMM
\end{enumerate}

\[ \gamma_i \sim N({\bf 0},{\bf Q})\] and independent from covariates.
Observe: we still assume normality for the random effects!

In addition, the mean \(\mu_{ij}=\text{E}(Y_{ij}\mid \gamma_i)\) is now
to be seen to be conditional on the random effect \(\gamma_i\). This
means that the specified response distribution is conditional on the
random effect: \(f(y_{ij}\mid \gamma_i)\).
\end{block}
\end{frame}

\begin{frame}
\begin{block}{GLMM model assumptions}
\phantomsection\label{glmm-model-assumptions}
\begin{block}{Distributional assumptions}
\phantomsection\label{distributional-assumptions}
(previously ``1. Random component from exponential family'')

\begin{itemize}
\tightlist
\item
  Given the random effects \(\gamma_i\) and the covariates
  \({\bf x}_{ij}\) and \({\bf u}_{ij}\)
\item
  the responses \(Y_{ij}\) are \emph{conditionally independent} and
\item
  the \emph{conditional distribution} \(f(y_{ij}\mid \gamma_i)\) belongs
  to and exponential family.
\end{itemize}
\end{block}
\end{block}
\end{frame}

\begin{frame}
\begin{block}{Structural assumptions}
\phantomsection\label{structural-assumptions}
(previously ``2. Systematic component'' and ``3. Link function'')

\begin{itemize}
\tightlist
\item
  The \emph{conditional mean} \(\mu_{ij}=\text{E}(Y_{ij}\mid \gamma_i)\)
  is linked to
\item
  the \emph{linear predictor}
  \(\eta_{ij}={\bf x}_{ij}^T {\boldsymbol \beta}+ {\bf u}_{ij} \gamma_i\)
\item
  through the link function \[\eta_{ij}=g(\mu_{ij})\]
\item
  or equivalently through the response function
  \[\mu_{ij}=h(\eta_{ij})\]
\item
  where \(h^{-1}=g\) (inverse functions).
\end{itemize}
\end{block}
\end{frame}

\begin{frame}
\begin{block}{Distributional assumptions for random effects}
\phantomsection\label{distributional-assumptions-for-random-effects}
\begin{itemize}
\tightlist
\item
  The random effects \(\gamma_i\), \(i=1,\ldots,m\) are independent and
  identically distributed \[ \gamma_i \sim N({\bf 0}, {\bf Q})\]
\item
  where the covariance matrix \({\bf Q}\) is a \((q+1)\times (q+1)\)
  positive definite.
\item
  An important special case is
  \({\bf Q}=diag(\tau_0^2, \tau_1^2,\ldots,\tau_q^2)\).
\end{itemize}

{Any questions? Notice that we still have normal distribution for the
random effects!}

\textbf{Remark}: for LMMs we used that
\(\varepsilon \sim N({\bf 0}, \sigma^2 {\bf I})\), but said that it was
possible to let the error terms be correlated
\(\varepsilon \sim N({\bf 0}, {\bf R})\) where \({\bf R}\) could have
non-zero off-diagonal elements. However, for GLMMs this is much more
complicated.
\end{block}
\end{frame}

\begin{frame}
\begin{block}{Alternative two-step formulation}
\phantomsection\label{alternative-two-step-formulation}
Consider a random variable \(Y_{ij}\) and a random effect \(\gamma_i\)
for cluster \(i\).

Alternative two step formulation:

\begin{enumerate}
\tightlist
\item
  \(\gamma_i \sim N({\bf 0},{\bf Q})\)
\item
  Given the realization of \(\gamma_i\), then the \(Y_{ij}\)s are
  independent and each have the distribution \(f(y_{ij}\mid \gamma_i)\)
  with mean \(\mu_{ij}=h(\eta_{ij})\) and linear predictor
  \(\eta_{ij}={\bf x}_{ij}^T {\boldsymbol \beta}+ {\bf u}^T_{ij} \gamma_i\).
\end{enumerate}

Remark: this conditional formulation is ``easy'' to write down, but to
do inference we need the marginal distribution of all \(Y_{ij}\)s
together to construct the likelihood. We will soon see that this is
hard.
\end{block}
\end{frame}

\begin{frame}{GLMM random intercept model}
\phantomsection\label{glmm-random-intercept-model}
\textbf{Distributional assumptions:} we will focus on binomial and
Poisson responses.

\textbf{Structural assumptions:} The random effect is added to the
linear predictor

\[\eta_{ij}={\bf x}_{ij}^T {\boldsymbol \beta}+ \gamma_{0i}\]

and we only consider canonical links, so logit for binomial and log for
Poisson.

\textbf{Distributional assumptions for random effects}

The \(\gamma_{0i}\), \(i=1,\ldots,m\) are independent and identically
distributed \[\gamma_{0i}\sim N(0,\tau_0^2)\]
\end{frame}

\begin{frame}
\begin{block}{Poisson random intercept model}
\phantomsection\label{poisson-random-intercept-model}
\textbf{Distributional assumptions:} The conditional distribution of the
response \(Y_{ij}\) is Poisson with mean \(\lambda_{ij}\) \[
f(y_{ij}\mid \gamma_i)=\frac{\lambda_{ij}^{y_{ij}}}{y_{ij}!}\exp(-\lambda_{ij}) \text{ for } y=0,1,2,\ldots
\] This is conditional on \(\lambda_{ij}\), which means that it is
really conditional on the linear predictor - since the mean is a
function of the linear predictor, and in the linear predictor we have
fixed and random effects. So, conditional on the values for the fixed
and random effects.

Then the observations \(Y_{ij}\) are \emph{conditionally} independent
for all \(i\) and \(j\) - but, not \emph{marginally} independent.
\end{block}
\end{frame}

\begin{frame}
\textbf{Structural assumptions:}
\[\eta_{ij}={\bf x}_{ij}^T {\boldsymbol \beta}+ \gamma_{0i}\]

\[\ln(\lambda_{ij})=\eta_{ij} \text{ or } \lambda_{ij}=\exp(\eta_{ij})\]

\textbf{Distributional assumptions for random effects}

The \(\gamma_{0i}\), \(i=1,\ldots,m\) are independent and identically
distributed \[\gamma_{0i}\sim N(0,\tau_0^2)\]
\end{frame}

\begin{frame}
\begin{block}{Overdispersion}
\phantomsection\label{overdispersion}
It is not necessary that \(n_i>1\) to use this model, and this model can
be used to take care of overdispersion (when all \(n_i=1\)).

For the Poisson case, this would be equivalent to using a Poisson
log-normal distribution
\end{block}
\end{frame}

\begin{frame}[fragile]
\begin{block}{Beach-example: Poisson GLMM with beach as random
intercept}
\phantomsection\label{beach-example-poisson-glmm-with-beach-as-random-intercept}
(we will talk about parameter estimation soon)

\begin{Shaded}
\begin{Highlighting}[]
\NormalTok{RI}\OtherTok{=}\FunctionTok{glmer}\NormalTok{(Richness}\SpecialCharTok{\textasciitilde{}}\NormalTok{NAP }\SpecialCharTok{+}\NormalTok{(}\DecValTok{1}\SpecialCharTok{|}\NormalTok{Beach),}\AttributeTok{data=}\NormalTok{RIKZ,}
         \AttributeTok{family=}\FunctionTok{poisson}\NormalTok{(}\AttributeTok{link=}\NormalTok{log))}
\FunctionTok{summary}\NormalTok{(RI)}\SpecialCharTok{$}\NormalTok{coefficients}
\end{Highlighting}
\end{Shaded}

\begin{verbatim}
##               Estimate Std. Error   z value     Pr(>|z|)
## (Intercept)  1.6623311 0.17372603  9.568693 1.082653e-21
## NAP         -0.5038922 0.07535222 -6.687158 2.275463e-11
\end{verbatim}

\begin{Shaded}
\begin{Highlighting}[]
\FunctionTok{summary}\NormalTok{(RI)}\SpecialCharTok{$}\NormalTok{varcor}
\end{Highlighting}
\end{Shaded}

\begin{verbatim}
##  Groups Name        Std.Dev.
##  Beach  (Intercept) 0.47428
\end{verbatim}
\end{block}
\end{frame}

\begin{frame}[fragile]
\begin{block}{Beach-example: Poisson GLMM with overdispersion}
\phantomsection\label{beach-example-poisson-glmm-with-overdispersion}
\texttt{Sample} is 1:nrow(RIKZ)

\begin{Shaded}
\begin{Highlighting}[]
\NormalTok{RIod}\OtherTok{=}\FunctionTok{glmer}\NormalTok{(Richness}\SpecialCharTok{\textasciitilde{}}\NormalTok{NAP }\SpecialCharTok{+}\NormalTok{(}\DecValTok{1}\SpecialCharTok{|}\NormalTok{Beach) }\SpecialCharTok{+}\NormalTok{(}\DecValTok{1}\SpecialCharTok{|}\NormalTok{Sample),  }
           \AttributeTok{data=}\NormalTok{RIKZ,}\AttributeTok{family=}\FunctionTok{poisson}\NormalTok{(}\AttributeTok{link=}\NormalTok{log))}
\FunctionTok{summary}\NormalTok{(RIod)}\SpecialCharTok{$}\NormalTok{coefficients}
\end{Highlighting}
\end{Shaded}

\begin{verbatim}
##              Estimate Std. Error   z value     Pr(>|z|)
## (Intercept)  1.636575 0.17475924  9.364739 7.623879e-21
## NAP         -0.551958 0.09671449 -5.707087 1.149261e-08
\end{verbatim}

\begin{Shaded}
\begin{Highlighting}[]
\FunctionTok{summary}\NormalTok{(RIod)}\SpecialCharTok{$}\NormalTok{varcor}
\end{Highlighting}
\end{Shaded}

\begin{verbatim}
##  Groups Name        Std.Dev.
##  Sample (Intercept) 0.26555 
##  Beach  (Intercept) 0.45822
\end{verbatim}
\end{block}
\end{frame}

\begin{frame}[fragile]
\tiny

\begin{Shaded}
\begin{Highlighting}[]
\FunctionTok{library}\NormalTok{(sjPlot)}
\FunctionTok{library}\NormalTok{(ggplot2)}

\FunctionTok{plot\_model}\NormalTok{(RI, }\AttributeTok{type =} \StringTok{"re"}\NormalTok{) }\SpecialCharTok{+} \FunctionTok{ylab}\NormalTok{(}\StringTok{"BLUP"}\NormalTok{) }\SpecialCharTok{+} \FunctionTok{xlab}\NormalTok{(}\StringTok{"Beaches"}\NormalTok{)}
\end{Highlighting}
\end{Shaded}

\pandocbounded{\includegraphics[keepaspectratio]{Module08GLMMPresentation_files/figure-beamer/unnamed-chunk-6-1.pdf}}
\end{frame}

\begin{frame}
\pandocbounded{\includegraphics[keepaspectratio]{Module08GLMMPresentation_files/figure-beamer/unnamed-chunk-7-1.pdf}}
\end{frame}

\begin{frame}
\pandocbounded{\includegraphics[keepaspectratio]{Module08GLMMPresentation_files/figure-beamer/unnamed-chunk-8-1.pdf}}
\end{frame}

\begin{frame}
\pandocbounded{\includegraphics[keepaspectratio]{Module08GLMMPresentation_files/figure-beamer/unnamed-chunk-9-1.pdf}}
\end{frame}

\begin{frame}
\pandocbounded{\includegraphics[keepaspectratio]{Module08GLMMPresentation_files/figure-beamer/unnamed-chunk-10-1.pdf}}

\normalsize
\end{frame}

\begin{frame}
\normalsize

{Comment on what you see. Would you expect that the solid and dashed
curve were identical?}
\end{frame}

\begin{frame}
\begin{block}{Binomial random intercept model}
\phantomsection\label{binomial-random-intercept-model}
We will study this model in the interactive session.

\begin{block}{Simulation example}
\phantomsection\label{simulation-example}
Aim: to show the difference between random intercept and random slope
logit models, and to compare the mean cluster curve with the population
curve (as for the Poisson model).

This example a slightly modified version (added line and changed to
ggplot) of an example in the Lecture notes by Magne Aldrin, Norwegian
Computing Centre (which is way the notation differs slightly from what
we have used).
\end{block}
\end{block}
\end{frame}

\begin{frame}
The model we use is
\[\begin{pmatrix} \gamma_{0i}\\ \gamma_{1i} \end{pmatrix} \sim N \left( \begin{pmatrix}0\\ 0 \end{pmatrix},\begin{pmatrix} 4 & 0 \\ 0 &1 \end{pmatrix} \right)\]
and then \[\eta_{ij} = \beta_0+\beta_1 x_{ij} + \gamma_{0i} \] for the
random intercept model and
\[\eta_{ij} = \beta_0+\beta_1 x_{ij} + \gamma_{0i}+\gamma_{1i}x_{ij} \]
for the random slope model where \(\beta_0=3\) and \(\beta_1=1\).
Logit-link is used so \(\eta_{ij}=\ln(\frac{\pi_{ij}}{1-\pi_{ij}})\).

Only data for the random intercept and slope are simulated, while
parameter estimates are considered known. The number of clusters is
\(m=1000\) and \(n_i = 201\) value of \(x\) from \(-10\) to \(10\) are
studied.
\end{frame}

\begin{frame}
Below 30 random lines are plotted together (in different colours). The
black dashed line is
\(\exp(\beta_0+\beta_i x)/(1+\exp(\beta_0+\beta_i x))\) (denoted ``not
mean'' in the plot) and the black solid line is the population mean
(really, the average of the \(m\) curves).
\end{frame}

\begin{frame}
\small

\pandocbounded{\includegraphics[keepaspectratio]{Module08GLMMPresentation_files/figure-beamer/unnamed-chunk-13-1.pdf}}

\normalsize

{Comment on what you see.}
\end{frame}

\begin{frame}{Conditional and marginal model}
\phantomsection\label{conditional-and-marginal-model}
To do parameter estimation in the LMM we started with the marginal model
for \(Y_{ij}\).

\[f(y_{ij})=\int_{\gamma_i} f(y_{ij}\mid \gamma_i) f(\gamma_i) d\gamma_i\]

where \[f(\gamma_i) \text{ is } N({\bf 0},{\bf Q})\] and
\(f(y_{ij}\mid \gamma_i)\) might be the binomial, Poisson, normal,
gamma, \ldots{} distribution - chosen as an exponential family.

It turns out that this integral can only be written out analytically in
special cases.

In the special case that \(f(y_{ij}\mid \gamma_i)\) is normal we saw in
Module LMM that the marginal distribution is also multivariate normal.
\end{frame}

\begin{frame}
Another special case is the log-linear Poisson random intercept where
all \(n_i=1\). See pages 393-394 in our textbook to see the solution.

For the random intercept Poisson GLMM with log-link the expected
marginal mean is (Agresti, 2015, page 309)
\[ \text{E}(Y_{ij})=\text{E}[\text{E}(Y_{ij}\mid \gamma_{0i})]=\exp({\bf x}_{ij}^T{\boldsymbol \beta}+\tau_0^2/2)\]
The solid line in the plot to the left above gives this curve, while the
dashed is \(\exp({\bf x}_{ij}^T{\boldsymbol \beta})\). In Agresti (2015,
page 310) the variance of the marginal distribution is given as
\[\text{Var}(Y_{ij})=\cdots = \text{E}(Y_{ij})+[\text{E}(Y_{ij})^2(\exp(\tau_0^2)-1)]\]
so we see that the variance is larger than that of a Poisson
distribution.

Remark: the conditional mean is
\(\exp({\bf x}_{ij}^T{\boldsymbol \beta}+\gamma_{0i})\), while the
marginal mean is \(\exp({\bf x}_{ij}^T{\boldsymbol \beta}+\tau_0^2/2)\).

(If the link function chosen had been the identity link then the
marginal mean would have been \({\bf x}_{ij}^T{\boldsymbol \beta}\).)
\end{frame}

\begin{frame}{Parameter estimation}
\phantomsection\label{parameter-estimation}
The most popular method for parameter estimation for regression is
maximum likelihood. The likelihood of the data is given by the marginal
distribution of all \(Y_{ij}\)s jointly. However, this can not be found
analytically (as we saw above), so we must resort to numerical methods
to evaluate the likelihood of a general GLMM.

The parameters we want to estimate are the fixed effects
\({\boldsymbol \beta}\)s, and the parameters in \({\bf Q}\) for the
random effects - denoted by \(\vartheta\).
\end{frame}

\begin{frame}
The contribution from observations in the \(i\)th cluster is

\[
\begin{aligned}
f({\bf y}_i\mid {\boldsymbol \beta},\vartheta) &= \int_{ {\boldsymbol \gamma}_i} f({\bf y}_i \mid {\boldsymbol \gamma}_i,{\boldsymbol \beta}) f({\boldsymbol \gamma}_i \mid {\bf Q}) d{\boldsymbol \gamma}_i
\\
&=\int_{ {\boldsymbol \gamma}_i} \prod_{j=1}^{n_i}f(y_{ij} \mid {\boldsymbol \gamma}_i,{\boldsymbol \beta}) f({\boldsymbol \gamma}_i \mid {\bf Q}) d{\boldsymbol \gamma}_i
\end{aligned}
\]

And the likelihood is then (since the clusters are independent)

\[
L({\boldsymbol \beta},\vartheta)=\prod_{i=1}^m f(\bf{y}_i\mid {{\boldsymbol \beta}}, \vartheta)
\]
\end{frame}

\begin{frame}
To arrive at parameter estimates we need to maximize the likelihood with
respect to \({\boldsymbol \beta}\) and \(\vartheta\), which is a
complicated numerical problem, because wee need to integtrate over the
\(\gamma_i\)'s.

Let
\(l({\bf y}_i, \boldsymbol{\gamma}_i \mid {\boldsymbol \theta}) = \log(({\bf y}_i, \boldsymbol{\gamma}_i \mid {\boldsymbol \beta},\vartheta))\)
be the \emph{joint log-likelihood of the data and the random effects}
(with \(\theta = (\beta, \vartheta)\)).

Observe that
\(l({\bf y}_i, \boldsymbol{\gamma}_i \mid {\boldsymbol \beta},\vartheta)\)
depends on the unknown random effects \(\gamma_i\) and the parameters in
\(\theta = (\beta, \vartheta)\) (here we have the fixed effects and the
parameters in \({\bf Q}\) for the random effects).
\end{frame}

\begin{frame}
The maximum likelihood estimate for \(\theta\) maximises

\[
L(\theta)=\int_{\boldsymbol{\gamma}} \exp(l({\bf y}_i, \boldsymbol{\gamma}_i \mid {\boldsymbol \theta}))d\gamma
\]

with respect to \(\theta\).

The random effects \(\gamma_i\) are now integrated out and the marginal
likelihood \(L(\theta)\) is the likelihood of the data as a function of
the parameters only.
\end{frame}

\begin{frame}[fragile]
\begin{block}{Approximations}
\phantomsection\label{approximations}
There are a few approaches to approximating the integral:

\begin{itemize}
\tightlist
\item
  The Laplace Approximation
\item
  Gauss-Hermite Quadrature
\item
  Simulation (usually MCMC: see TMA4300 for more)
\item
  INLA (Bayesian, much more advance. Involves nested Laplace
  approximations)
\end{itemize}

\texttt{glmer} function in \texttt{lme4} uses the Laplace approximation
as a default, but can also use Gauss-Hermite Quadrature. Bayesian
methods abound: MCMC is slow, but more flexible.
\end{block}
\end{frame}

\begin{frame}
\begin{block}{The Laplace Approximation}
\phantomsection\label{the-laplace-approximation}
The Laplace Approximation assumes that
\(l({\bf y}_i, \boldsymbol{\gamma} \mid {\boldsymbol \theta})\) is
quadratic in \(\boldsymbol{\gamma}\): this is equivalent to either

\begin{itemize}
\tightlist
\item
  assuming a Gaussian distribution for the sampling distribution of each
  \(\gamma_i\), or
\item
  using a single quadrature point in Gauss-Hermite Quadrature
\end{itemize}
\end{block}
\end{frame}

\begin{frame}
\begin{block}{The Approximation}
\phantomsection\label{the-approximation}
Let \(g(\gamma)=\ln(f(y\mid \gamma)f(\gamma))\), and look at the second
order Taylor expansion around \(\hat{u}(\theta)\):

\[
\begin{aligned}
g(\gamma_i) &\approx \tilde{g}(\gamma_i)=g(\hat{\gamma_i})+(\gamma_i-\hat{\gamma_i})g'(\hat{\gamma_i})+
\frac{1}{2}(\gamma_i-\hat{\gamma_i})^2 g''(\hat{\gamma_i}) \\
&=g(\hat{\gamma_i})-\frac{1}{2}(\gamma_i-\hat{\gamma_i})^2 (-g''(\hat{\gamma_i}))
\end{aligned}
\]

This means that \(\exp(\tilde{g}(\gamma_i))\) is proportional to the
normal density with mean \(\mu=\hat{\gamma_i}\) and variance
\(\sigma^2=-1/g''(\hat{\gamma_i})\).

Putting this back into our integral:

\[
\begin{aligned}
L(\theta) &= \int_{\boldsymbol{\gamma}} \exp(l({\bf y}_i, \boldsymbol{\gamma}_i \mid {\boldsymbol \theta}))d\boldsymbol{\gamma} \approx\int_{\boldsymbol{\gamma}} \exp(\tilde{l}({\bf y}_i, \boldsymbol{\gamma}_i \mid {\boldsymbol \theta}))d\boldsymbol{\gamma} \\
&=\exp(l(\bf{y}_i, \hat{ {\boldsymbol \gamma}} \mid \boldsymbol{\theta}))
\int_{\boldsymbol{\gamma}} \exp(-\frac{1}{\sigma^2} (\boldsymbol{\gamma}-\hat{\boldsymbol{\gamma}})^2) d\boldsymbol{\gamma} \\
&=\exp(l({\bf y}_i, \hat{\boldsymbol{\gamma}} \mid {\boldsymbol \theta}))\sqrt{-2\pi/g''(\hat{\gamma_i})}
\end{aligned}
\]
\end{block}
\end{frame}

\begin{frame}
Now we can plus this into our full likelihood, pausing only to note that
\(H(\theta)\) is the Hessian of
\(l({\bf y}_i, \boldsymbol{\gamma}_i \mid {\boldsymbol \theta})\)
evaluated at \(\hat{\boldsymbol{\gamma}}(\theta)\), so contains
\(g''(\hat{\gamma}_i)\) terms

Let \(\hat{\boldsymbol{\gamma}}(\theta)\) be the maximum of
\(l({\bf y}_i, \boldsymbol{\gamma}_i \mid {\boldsymbol \theta})\) w.r.t
\(\boldsymbol{\gamma}\).

The Laplace approximation for the marginal likelihood \(L(\theta)\) is
then

\[
L^{*}(\theta)=(2\pi)^{n/2} \text{det}(H(\theta))^{-1/2} \exp(f(\hat{u},\theta))
\]
\end{frame}

\begin{frame}
\begin{block}{Gauss-Hermite Quadrature}
\phantomsection\label{gauss-hermite-quadrature}
If the dimension of \({\boldsymbol \gamma}_i\) is small, for example our
random intercept model, then the Gauss-Hermite quadrature may be used to
approximate the integral above. This is because the form

\[\int_{-\infty}^{\infty} h(\gamma)\exp(-\gamma^2)d\gamma\]

can be written as a sum of weights and quadrature points that are roots
of Hermite polynomials. This can then be maximized using Newton-Raphson
methods.

When the dimension of \({\boldsymbol \gamma}\) gets large, this gets too
slow, so other methods are usually better.

In essence, this is a better approximation than the Laplace
Approximation, but there is a computational cost that might not be worth
paying.
\end{block}
\end{frame}

\begin{frame}{What have we not covered?}
\phantomsection\label{what-have-we-not-covered}
\begin{itemize}
\tightlist
\item
  We have skipped a lot of technicalities about the parameter
  estimation.
\item
  How to predict the random effects \(\gamma_i\): ``the conditional
  modes'' of the random effects.
\item
  What is fitted values, and how are residuals calculated.
\item
  How to test hypotheses?
\item
  AIC to compare models (very similar to LMM).
\item
  And, surely, much more.
\end{itemize}
\end{frame}

\begin{frame}{Summing up - what have we learned about the GLMM?}
\phantomsection\label{summing-up---what-have-we-learned-about-the-glmm}
\begin{itemize}
\tightlist
\item
  The GLMM can be formulated using three ingredients:

  \begin{itemize}
  \tightlist
  \item
    distributional assumption: \(f(y_{ij}\mid \gamma_i)\) from
    exponential family
  \item
    structural assumptions: linear predictor
    \(\eta_{ij}={\bf x}_{ij}^T {\boldsymbol \beta}+{\bf u}_{ij}\gamma_i\)
    and link function \(\eta_{ij}=g(\mu_{ij})\) where
    \(\mu_{ij}=\text{E}(Y_{ij}\mid \gamma_i)\).
  \item
    distributional assumptions for the random effects:
    \(\gamma_i\sim N({\bf 0},{\bf Q})\).
  \end{itemize}
\item
  The GLMM likelihood function is expressed as an integral with respect
  to the random effects and does (in general) not have a closed form
  solution.
\item
  Numerical approximation methods need to be used to find parameter
  estimates, and one possibility is to use the Laplace approximation,
  but many competing method exists.
\end{itemize}
\end{frame}

\begin{frame}[fragile]
Three \texttt{R} packages that can be investigated are

\begin{itemize}
\tightlist
\item
  \texttt{lme4}(function \texttt{glmer}),
\item
  \texttt{glmmTMB} (template model builder),
\item
  and the NTNU-flagship \texttt{inla}.
\end{itemize}

How to use these three packages on a simulated data set (binary data,
logit link, random intercept and slope) is shown in the end of the
module page.

For those who want to use MCMC, the \texttt{MCMCglmm} package can be
used, or you can use general MCMC programmes, such as JAGS, Nimble and
STAN (the only real advantage of these nowadays is if you want to go
beyond GLMMs)
\end{frame}

\begin{frame}[fragile]{R packages}
\phantomsection\label{r-packages}
\begin{Shaded}
\begin{Highlighting}[]
\CommentTok{\# install.packages(\textquotesingle{}arm\textquotesingle{})}
\FunctionTok{install.packages}\NormalTok{(}\StringTok{"reshape2"}\NormalTok{)}
\FunctionTok{install.packages}\NormalTok{(}\StringTok{"sp"}\NormalTok{)}
\FunctionTok{install.packages}\NormalTok{(}\StringTok{"glmmTBM"}\NormalTok{)}
\FunctionTok{install.packages}\NormalTok{(}\StringTok{"INLA"}\NormalTok{, }\AttributeTok{repos =} \StringTok{"https://inla.r{-}inla{-}download.org/R/stable"}\NormalTok{, }\AttributeTok{dep =} \ConstantTok{TRUE}\NormalTok{)}
\FunctionTok{install.packages}\NormalTok{(}\StringTok{"lme4"}\NormalTok{)}
\FunctionTok{install.packages}\NormalTok{(}\StringTok{"devtools"}\NormalTok{)}
\FunctionTok{library}\NormalTok{(devtools)}
\CommentTok{\# install\_github(\textquotesingle{}romunov/AED\textquotesingle{})}
\FunctionTok{install.packages}\NormalTok{(}\StringTok{"sjPlot"}\NormalTok{)}
\FunctionTok{install.packages}\NormalTok{(}\StringTok{"sjmisc"}\NormalTok{)}
\end{Highlighting}
\end{Shaded}
\end{frame}

\begin{frame}{Further reading}
\phantomsection\label{further-reading}
\begin{itemize}
\tightlist
\item
  \href{http://avesbiodiv.mncn.csic.es/estadistica/curso2011/regm26.pdf}{Bolker
  et al.~(2008):\emph{Generalized linear mixed models: a practical guide
  for ecology and evolution}}
\item
  Zuur et al.~(2009): ``Mixed Effects Models and Extensions in Ecology
  with R'', chapter 13 (pages 323-341). Available as free ebook from
  Springer for NTNU students. More explanations and less mathematics
  than Fahrmeir et al (2013), more focus on understanding.
  \href{https://link.springer.com/chapter/10.1007/978-0-387-87458-6_13}{Link
  to ebook Chapter 13}
\item
  Agresti (2015): Foundations of linear and generalized linear models,
  Chapter 9.
\end{itemize}
\end{frame}

\end{document}
